% !TEX root = ../cubic-moduli-descent.tex
\section{局所一ブロック型の構造対象付き圏モジュライ}
\label{sec:nfm-local-moduli}

本節では、固定表示模型間の全同値比較を対象の効果化へ進める。
扱う族の範囲を、良い一ブロック模型を fpqc 局所的に持つ族に限定する。
この局所表示条件を満たすこと自体を、任意の非可換三次多様体に対して
証明するわけではない。
一方、この範囲では次数一 shift を標識として要求せず、
構造対象を持つ Grothendieck 圏の族を代数的 stack として分類する。

次数スカラーを除く商と、その上の次数別に重みを持つ twisted algebra は
Behrend--Noohi\cite[Section~2.3]{NFMBehrendNoohi}の一般的構成に既に現れる。
同論文の偏極付き Grothendieck 圏は自己同値もデータに含む
\cite[Section~3.1.2]{NFMBehrendNoohi}。
ここで必要な追加点は、定理\ref{thm:nfm-fixed-model-comparison}により
\emph{全ての} $\mathcal O$-保存同値を表示の同型へ戻すことと、
twisted algebra から大域の $\mathcal O$-marked 圏を実際に構成することである。
後半では、一ブロックの中心階数と active 階数から
持上げ障害の正確な位数を調べる。

\subsection{係数変更を備えた圏族}

$R$ を可換 $\mathbf C$-代数とする。
$R$-線型 Grothendieck Abel 圏 $\mathcal C$ と可換 $R$-代数 $T$ に対し、
\[
 \mathcal C_T=\operatorname{Mod}_T(\mathcal C)
\]
は $T$-作用付き対象の圏を表す。
係数延長は copower $X\mapsto T\otimes_RX$ であり、
反復係数変更には標準の結合同型を用いる。
したがって $(\mathcal C,\mathcal O)$ は、各 affine $R$-scheme 上の
圏と構造対象 $(\mathcal C_T,T\otimes_R\mathcal O)$ を定める。
この係数変更構造を持つ族だけを以下で扱い、
幾何学的ファイバーの圏を単に並べたものとは区別する。

\begin{lemma}[Grothendieck 圏内の忠実平坦降下]
\label{lem:nfm-grothendieck-fpqc}
$\mathcal C$ を $R$-線型 Grothendieck Abel 圏とする。
任意の忠実平坦 $R\to T$ に対して、
\[
 T\otimes_R-:\mathcal C\longrightarrow\mathcal C_T
\]
は完全かつ保守的であり、$\mathcal C$ は $\mathcal C_T$ の
通常の fpqc descent datum の圏に同値である。
また、二つのこのような係数圏族の間の線型完全同値、自然同型、
および指定対象上の同定は fpqc 降下を満たす。
\end{lemma}

\begin{proof}
加群 $M$ による copower $M\otimes_R X$ は、自由加群による表示を用いて
余核として構成でき、$M$ と $X$ のそれぞれについて colimit を保つ。
Lazard の定理\cite[Theorem~10.81.4, Tag~058G]{NFMStacksLazard}により、
平坦 $T$ は有限自由加群の filtered colimit である。
Grothendieck 圏の filtered colimit は完全なので、$T\otimes_R-$ は完全である。

保守性のために、任意の $X$ に対して $X\to T\otimes_RX$ が単射であることを示す。
忠実平坦性により $R\to T$ は純単射で、$Q=T/R$ は平坦である。
実際、任意の $R$-加群 $M$ に対する $M\to T\otimes_RM$ は、
$T$ を tensor した後には積写像による左逆を持つため単射であり、
$\operatorname{Tor}_1^R(Q,M)=0$ が従う。
$Q=\varinjlim Q_i$ を有限自由加群による表示とし、
$T_i=T\times_QQ_i$ と置くと
\[
 0\longrightarrow R\longrightarrow T_i\longrightarrow Q_i\longrightarrow0
\]
は分裂完全で、$\varinjlim T_i=T$ である。
これらに $-\otimes_RX$ を施し、完全な filtered colimit を取ると
$X\to T\otimes_RX$ は単射となる。
完全性と合わせて保守性を得る。

対象と射の降下は、補題\ref{lem:nfm-tails-fpqc}の等化子による証明を適用する。
必要な入力は、今証明した完全性・保守性と Abel 圏内の核だけである。
具体的に、descent datum を持つ $Y\in\mathcal C_T$ の coaction を
$\theta:UY\to T\otimes_RUY$ と書けば、
\[
 X=\ker\bigl(\theta-\eta:UY\rightrightarrows T\otimes_RUY\bigr)
\]
が降下対象である。cocycle 恒等式と $T$-作用の counit 恒等式から
$T\otimes_RX\simeq Y$ を得る。
同じ等化子は射にも適用できる。

関手と自然同型の降下には命題\ref{prop:nfm-equivalence-descent}の証明を用いる。
局所同値を降下対象ごとに適用して関手を作り、逆同値も降下させる。
得られた関手の係数変更が元の関手と一致することは、自由 $T$-作用付き対象と
標準の reflexive coequalizer
\[
 T\otimes_RT\otimes_RUY\rightrightarrows T\otimes_RUY\longrightarrow Y
\]
上で確認する。同値は colimit を保つので、この確認で全対象を扱える。
指定対象上の同定は降下対象間の射として降下する。
\end{proof}

以後 $r=n-1\ge2$, $N=n+2=r+3$ とする。
第\ref{sec:nfm-graded-stack}節の $T^{\mathrm{fin}}$ を、
必要ならさらに $P$-不変な非空開 $T^{\mathrm{good}}$ へ制限する。
この開上の一ブロック模型は、普遍中心条件、中心上の平坦性、
次数別有限局所自由性、および
定理\ref{thm:nfm-fixed-model-comparison}に必要な相対双対性を
全ての係数変更後に満たすものとする。
相対 PBW 分裂と純 active Koszul CY 入力は
第\ref{sec:nfm-relative-duality}節の仮定である。
これらは本節で追加の分類定理として証明するものではない。
中心三次の torsor をこの開へ制限して $J^{\mathrm{good}}$ と書き、
\[
 \mathfrak C_n^{\mathrm{good}}=[J^{\mathrm{good}}/\overline P]
\]
と置く。以下の議論はこの stack の任意の非空開部分にも制限できる。

\begin{definition}[局所一ブロック型の圏族]
\label{def:nfm-local-category}
$\mathfrak M_n^{\mathrm{loc}}(R)$ の対象は、
$R$-線型 Grothendieck Abel 圏と指定対象の対 $(\mathcal C,\mathcal O)$ であって、
ある忠実平坦 $R$-代数 $T$ と $j\in J^{\mathrm{good}}(T)$ に対し
\[
 (\mathcal C_T,\mathcal O_T)
 \simeq (\operatorname{QGr} A_j,\mathcal O_{A_j})
\]
となるものとする。
射は $R$-線型完全同値 $F$ と同定
$\alpha:F(\mathcal O)\xrightarrow{\sim}\mathcal O'$ の対、
二射は $\alpha$ を保つ自然同型とする。
圏族の係数変更は上で定めた $\operatorname{Mod}_T$ によるものとし、
次数一 shift、次数一生成元の束、および局所表示は対象の標識に含めない。
\end{definition}

\subsection{twisted graded module による対象の効果化}

$x\in\mathfrak C_n^{\mathrm{good}}(R)$ とし、
対応する $\overline P$-torsor を $Q\to\operatorname{Spec}R$ とする。
中心拡大
\[
 1\longrightarrow\mathbf G_m\longrightarrow P
 \longrightarrow\overline P\longrightarrow1
\]
に関する $Q$ の持上げ gerbe を $\mathcal G_x$ と書く。
これは $x$ による $[J^{\mathrm{good}}/P]$ の引戻しである。
普遍次数付き代数は $\mathcal G_x$ 上の代数
$\mathcal A=\bigoplus_{j\ge0}\mathcal A_j$ を与え、
inertia の $\lambda\in\mathbf G_m$ は $\mathcal A_j$ 上で $\lambda^j$ として作用する。
この規約を「重み $j$」と呼ぶ。

\begin{lemma}[重み付き tails の大域模型]
\label{lem:nfm-weighted-tails}
$\operatorname{Gr}^{\mathrm{wt}}_{\mathcal G_x}(\mathcal A)$ を、
各 $M_j$ が重み $j$ の準連接層であるような次数付き右
$\mathcal A$-加群 $M=\bigoplus_{j\in\mathbf Z}M_j$ の圏とする。
fpqc 局所的に通常の graded torsion となる加群全体を
$\operatorname{Tors}^{\mathrm{wt}}$ と書けば、
\[
 \mathcal C_x=
 \operatorname{Gr}^{\mathrm{wt}}_{\mathcal G_x}(\mathcal A)
 /\operatorname{Tors}^{\mathrm{wt}},\qquad
 \mathcal O_x=\pi\mathcal A
\]
は $R$-線型 Grothendieck Abel 圏とその指定対象を定める。
任意の $R$-代数 $T$ について標準同値
\[
 (\mathcal C_x)_T\simeq\mathcal C_{x_T}
\]
があり、構造対象と反復係数変更を保つ。
gerbe を中立化する fpqc 被覆上では、これは通常の
$(\operatorname{QGr}A,\mathcal O_A)$ に同値である。
\end{lemma}

\begin{proof}
代数的 stack 上の準連接層の圏は Grothendieck 圏である
\cite[Proposition~96.15.2, Tag~0781]{NFMStacksQCohGrothendieck}。
$\mathbf G_m$-gerbe 上の準連接層には重み分解があり、
各重みへの射影は完全で colimit を保つ。
この主張は gerbe を中立化した後の $\mathbf G_m$-表現の重み分解で確認でき、
fpqc 降下により大域化する。
従って各重みの圏、さらにその全整数にわたる積も Grothendieck 圏である。
ここで generator は各重みへ射影した generator の直和から得られる。

重み付き次数付き層から自由右 $\mathcal A$-加群を作る関手は
\[
 (E_i)_i\longmapsto
 \left(\bigoplus_{i\le j}E_i\otimes\mathcal A_{j-i}\right)_j
\]
である。加群の核、余核、filtered colimit は各次数で計算でき、
この自由加群関手を generator に適用すると加群圏の generator が得られる。
したがって $\operatorname{Gr}^{\mathrm{wt}}(\mathcal A)$ は Grothendieck 圏である。
次数一の有限生成性により torsion は部分対象、商、拡大、直和で閉じる。
このことは fpqc 局所で通常の次数付き加群の証明に帰着する。
従ってこれは localizing 部分圏で、表示した Gabriel 商も Grothendieck 圏である。

次に係数変更を調べる。
$\mathcal G_{x_T}\to\mathcal G_x$ は affine base change であるから、
その上の準連接層は $\mathcal G_x$ 上の $T$-作用付き準連接層に等しい。
この同一視は affine atlas 上の加群の同一視と、その fpqc 降下から従い、
inertia の重みおよび $\mathcal A$-作用を保つ。
従って
\[
 \operatorname{Gr}^{\mathrm{wt}}_{\mathcal G_{x_T}}(\mathcal A_T)
 \simeq\operatorname{Mod}_T
 \bigl(\operatorname{Gr}^{\mathrm{wt}}_{\mathcal G_x}(\mathcal A)\bigr).
\]
Gabriel 商の完全左随伴と完全忠実右随伴に $T$-作用を持たせる。
補題\ref{lem:nfm-coefficient-objects}と同様に、その核はちょうど
$\mathcal A_T$-torsion である。
これで任意の、非平坦なものも含む係数変更との整合を得る。

最後に $\mathcal G_x$ を中立化すると、各重み $j$ の準連接層は通常の
$R$-加群に一意に対応し、$\mathcal A$ は通常の次数付き代数になる。
重み付き加群・torsion・構造対象はそれぞれ通常のものへ移る。
\end{proof}

\begin{theorem}[局所一ブロック型の $\mathcal O$-marked モジュライ]
\label{thm:nfm-local-moduli}
上記の良い表示の入力の下で、
\[
 \mathfrak C_n^{\mathrm{good}}
 \xrightarrow{\ \sim\ }
 \mathfrak M_n^{\mathrm{loc}},\qquad
 x\longmapsto(\mathcal C_x,\mathcal O_x)
\]
は fpqc stack の同値である。
右辺の $2$-groupoid は通常の groupoid と同値であり、
指定対象上で恒等となる自然自己同型は全て恒等である。
特に $\mathfrak M_n^{\mathrm{loc}}$ は滑らかな既約有限型 DM stack であり、
\[
 \dim\mathfrak M_n^{\mathrm{loc}}
 =\binom{n+4}{3}-(n+2)^2+3n-2.
\]
この結論は、任意の局所模型の descent datum に対し、
大域の Grothendieck Abel 圏とその構造対象が有効に存在することを含む。
\end{theorem}

\begin{proof}
まず射を比較する。二つの表示 $A=B/(w)$ と $A'=B'/(w')$ の間の
次数付き同型は、次数二の関係空間から $B\simeq B'$ を誘導する。
実際、三次関係は次数二の関係を変えず、$B$ はその二次関係で定義される。
得られた同型は三次の核直線を互いに移すので、
命題\ref{prop:nfm-graded-presentation}により $J^{\mathrm{good}}$ の $P$-作用の射に一致する。
定理\ref{thm:nfm-fixed-model-comparison}により、次数スカラーを除いたこの射の層は
全ての $\mathcal O$-保存同値の層に一致する。
したがって主張する射は fpqc 局所で完全忠実である。
補題\ref{lem:nfm-grothendieck-fpqc}により、この完全忠実性は大域化する。

局所模型の同値の自然自己同型はスカラーであり、
その $\mathcal O$ 上の成分を恒等に指定するとスカラーは一となる
\textup{(}補題\ref{lem:nfm-natural-scalars}\textup{)}。
一般の対象でも、自然自己同型を良い被覆へ延長してこの結論を適用し、
忠実平坦降下で恒等性を確認できる。
従って各同値の groupoid は離散 groupoid に同値である。
ここで構造対象の標識を忘れた二自己同型の $\mathbf G_m$ と、
表示の次数スカラーを除いた通常の stack とを混同しない。

次に $(\mathcal C,\mathcal O)\in\mathfrak M_n^{\mathrm{loc}}(R)$ を取り、
定義の被覆上で $(\mathcal C_T,\mathcal O_T)$ を模型 $A_j$ と同定する。
重なりで得られる marked equivalence の類は、前段により
商 stack の射に一意に対応する。
元の係数圏族の結合同型は cocycle を与え、
marked natural automorphism の自明性により三重重なりでの曖昧さはない。
商 stack の対象降下から $x\in[J^{\mathrm{good}}/\overline P](R)$ を得る。
補題\ref{lem:nfm-weighted-tails}で構成した $(\mathcal C_x,\mathcal O_x)$ は
同じ局所模型と同じ同値の cocycle を持つ。
両者間の局所同値は補題\ref{lem:nfm-grothendieck-fpqc}で降下するので、
$(\mathcal C,\mathcal O)\simeq(\mathcal C_x,\mathcal O_x)$ である。

同じ議論は、最初から大域の $(\mathcal C,\mathcal O)$ が与えられていない
局所対象と marked equivalence の descent datum にも適用できる。
商 stack 上の対象を先に降下し、重み付き tails を作ることで、
大域の Grothendieck 圏を得る。
したがって対象の効果化を定義に含めて仮定したのではない。
滑らかさ、既約性、有限型性、DM 性と次元は
定理\ref{thm:nfm-central-torsor}から従う。
\end{proof}

\begin{remark}[証明した範囲]
局所的に良い一ブロック表示を持つという条件の下で、
標識は $(\mathcal C,\mathcal O)$ だけで十分である。
ただし、この条件が全ての内在的モジュライ成分の一般点で成立すること、
この stack が未構成の全モジュライの開かつ閉な成分をなすこと、
任意の Abel 圏の変形が局所表示条件を保つことは、この定理には含まれない。
これらには独立の局所表示・変形・成分の定理が必要である。
\end{remark}

\subsection{大域次数付き表示の gerbe と三捩れ障害}

\begin{proposition}[表示の持上げ gerbe]
\label{prop:nfm-presentation-gerbe}
写像
\[
 \rho_n:[J^{\mathrm{good}}/P]\longrightarrow
 \mathfrak M_n^{\mathrm{loc}}
\]
は $\mathbf G_m$-gerbe である。
その class を $\beta$ とし、$e=\gcd(r,3)=\gcd(n-1,3)$ と置くと
\[
 r\beta=3\beta=0,\qquad e\beta=0.
\]
任意の affine $S$ と $x\in\mathfrak M_n^{\mathrm{loc}}(S)$ に対し、
$\beta_x=0$ であることと、$x$ が同じ良い一ブロック型の
大域次数付き代数の tails として表示されることは同値である。
ここで次数一は有限局所自由でよく、自由な基底や中心三次直線の生成元は要求しない。

特に $n\not\equiv1\pmod3$ ならこの gerbe は全体で中立である。
$ar+3b=1$ となる整数 $a,b$ を一つ固定すると、
\[
 \chi(g)=\det(g|_U)^a\det(g|_{V/U})^b
\]
がその中立化を与える。
\end{proposition}

\begin{proof}
次数スカラーは $J^{\mathrm{good}}$ に自明に作用するため、
$P$ から $\overline P$ への商の fibre は
その torsor の $P$-持上げの gerbe である。
$[J^{\mathrm{good}}/P]$ 上には、重み一の階数 $r$ の束 $\mathcal U$ と
階数三の束 $\mathcal V/\mathcal U$ がある。
その determinant はそれぞれ重み $r$ と三の可逆層である。
gerbe 上の重み $m$ の可逆層は $m\beta=0$ を与える。
これは局所表示の三重重なりの scalar cocycle $c_{ijk}$ に対し、
可逆層の遷移写像が $c_{ijk}^{m}$ を coboundary として消すことから分かる。
従って表示した annihilator を得る。

$\beta_x=0$ は持上げ gerbe の中立性、すなわち $P$-torsor への持上げの存在と同値である。
持上げを取れば $P$-同変な普遍代数が $S$ 上へ降下し、
次数一は有限局所自由、中心一次部分は階数 $r$ の部分束になる。
三次関係は中心三次の可逆部分束で生成される。
逆にこの種の大域次数付き表示の、中心一次部分を保つ frame torsor は
$P$-torsor となり、$x$ を $[J^{\mathrm{good}}/P]$ へ持ち上げる。

$e=1$ のとき、$\chi(\lambda\operatorname{id}_V)=\lambda$ である。
写像 $p\mapsto\chi(p)^{-1}p$ は $\overline P$ の $\ker\chi$ への同型を与え、
\[
 P\simeq\mathbf G_m\times\ker\chi,\qquad
 \overline P\simeq\ker\chi
\]
と中心拡大を分裂する。
従って全ての $\overline P$-torsor はこの分裂により持ち上がる。
同じことは、gerbe 上の普遍次数 $j$ の代数束を
$\mathcal A_j\otimes\chi^{-j}$ に置き換えると重みが零になることからも分かる。
これは同じ局所代数を表すが、選んだ整数 $a,b$ に依存する中立化である。
\end{proof}

\begin{remark}[対象の障害と同値の障害]
$\beta_x\in H^2_{\mathrm{fpqc}}(S,\mathbf G_m)$ は
\emph{対象}が大域次数付き表示を持つかどうかの障害である。
一つの持上げが存在すれば、他の持上げとの差は基底の可逆層 $L$ であり、
代数は $\mathcal A_j\mapsto\mathcal A_j\otimes L^{\otimes j}$ と変わる。
一方、命題\ref{prop:nfm-lifting-line}の $L_F$ は、
\emph{既に二つの大域模型を固定した後}の指定同値の持上げ障害である。
後者の $H^1$ の現象から前者の $H^2$ の効果化を推測してはならない。
また $\beta_x=0$ は次数一が自由であること、中心三次直線が自明であること、
あるいは次数一 shift の標準的な選択を意味しない。
\end{remark}

\begin{theorem}[三捩れ障害はどの非空開でも消えない]
\label{thm:nfm-brauer-dense-open}
$n\equiv1\pmod3$、すなわち $r=3m$ とする。
任意の非空開部分 stack
$\mathfrak U\subset\mathfrak M_n^{\mathrm{loc}}$ に対し、
\[
 K=\mathbf C(s,t),\qquad x\in\mathfrak U(K)
\]
で $\beta_x\ne0$ となるものが存在する。
従って $\beta|_{\mathfrak U}$ は位数がちょうど三であり、
$\rho_n$ はどの非空開部分上でも中立化しない。
この主張は非空開全体についてのもので、特殊次元や境界点の例に依存しない。
\end{theorem}

\begin{proof}
まず $K$ 上の非自明な degree 三の中心 division algebra を固定する。
$\zeta$ を原始三乗根とし、
\[
 D=(s,t)_3
 =K\langle a,b\rangle/
 (a^3-s,\ b^3-t,\ ba-\zeta ab)
\]
と置く。$a^i b^j$ $(0\le i,j<3)$ は基底であり、中心は $K$ である。
division 性は $K'=\mathbf C(s)((t))$ へ延長して確認できる。
$L=\mathbf C(s)(a)$、$a^3=s$ とすると $[L:\mathbf C(s)]=3$ で、
$D\otimes_KK'$ は
\[
 f_0+f_1b+f_2b^2,\qquad f_i\in L((t)),\qquad
 bf=\sigma(f)b,\quad\sigma(a)=\zeta a
\]
と書ける。非零元について
\[
 v(f_0+f_1b+f_2b^2)
 =\min_i\bigl(v_t(f_i)+i/3\bigr)
\]
と定める。
三つの候補値は小数部分が異なるので最小項は一意であり、
積の最小項も一意で消えない。従って積の valuation は和に等しく、
この有限次元代数に零因子はない。
有限次元代数の非零元による乗法は単射なら全射なので、これは division algebra である。
よって $D$ 自身も division algebra であり、その Brauer class は非零である。

$U\simeq\mathbf C^3\otimes\mathbf C^m$、$V/U\simeq\mathbf C^3$ を選び、
\[
 \operatorname{GL}_3\longrightarrow P,\qquad
 g\longmapsto\operatorname{diag}(g^{\oplus m},g)
\]
を考える。中心スカラーは $V$ 全体の同じスカラーへ移るから、
$\operatorname{PGL}_3\to\overline P$ を誘導する。
この写像に沿って $P\to\overline P$ を引き戻した中心拡大は
$\operatorname{GL}_3\to\operatorname{PGL}_3$ そのものである。
$D$ に対応する $\operatorname{PGL}_3$-torsor
\cite[Tag~03R6]{NFMStacksBrauerForms}をこの写像で押し出して
$\overline P$-torsor $Q$ を得ると、その持上げ障害は $[D]$ である。
Brauer class の符号規約を逆に取っても、以下の非零性には影響しない。

$\mathfrak U$ の逆像を $J'\subset J^{\mathrm{good}}$ とする。
これは非空な $\overline P$-不変開である。
表示 potential の空間
\[
 W_F=\operatorname{Sym}^3V/\operatorname{Sym}^3U
\]
には $\det(V/U)^{-1}$ で捩じった作用を用いており、
次数スカラーの重みは零である。
従って $W_F$ は $\overline P$ の線型表現であり、
$Q\times^{\overline P}W_F$ は $K$ 上の通常の有限次元ベクトル空間となる。
$T^{\mathrm{good}}_Q$ はその非空開である。
さらに $J^{\mathrm{good}}_Q\to T^{\mathrm{good}}_Q$ は
ベクトル束の加法群による torsor であり、特に滑らかである。

非空開 $J'_Q$ の $T^{\mathrm{good}}_Q$ における像は非空開である。
$K$ は無限体なので、この像には $K$-有理点 $z$ がある。
$z$ 上の fibre $J^{\mathrm{good}}_{Q,z}$ は有限次元ベクトル空間の加法群の torsor であり、
体上では自明である。
ここでは characteristic zero のため、有限 Galois 拡大上の一点の
Galois 軌道の affine barycenter を取ることでも有理点を構成できる。
従ってこの fibre は affine 空間であり、その非空開 $J'_{Q,z}$ も $K$-有理点を持つ。
この一点は $Q\to J'$ の同変写像、すなわち
$x\in[J'/\overline P](K)=\mathfrak U(K)$ を与える。
$x$ の持上げ gerbe の障害は固定した $Q$ の障害 $[D]\ne0$ である。
命題\ref{prop:nfm-presentation-gerbe}の $3\beta=0$ と合わせて、位数はちょうど三となる。
\end{proof}

\begin{remark}[一般点と幾何学的一般点]
この非中立性は、非空な稠密開に制限するだけでは
大域次数付き表示の障害を取り除けないことを意味する。
代数閉体上の個々の幾何学的点では gerbe は中立であり、矛盾はない。
定理は商 stack の幾何学的一般点に非零の体の Brauer class が残るという意味ではなく、
\emph{全ての非空開部分 stack 上で普遍 gerbe が非中立である}という主張である。
\end{remark}

\begin{remark}[Poisson 標的の偏極に対する制約]
後に構成する射がこの表示から得られる ambient 射影空間束を保つとする。
その ambient は $\overline P\to\operatorname{PGL}(V)$ による
Brauer--Severi 束であり、符号を除いて class は $\beta$ である。
もし標的を、実際の相対可逆層 $L$ を持ち、$p_*L$ が階数 $N$ で
基底変更と両立し、$L$ の全切断がその ambient 埋込みを回復する三次族に限定すれば、
ambient は $\mathbf P((p_*L)^\vee)$ となり Brauer class は零でなければならない。
従って $n\equiv1\pmod3$ では、その種の標的への ambient と両立する射を
この stack の非空開全体に作るには、gerbe の中立化という追加データが必要である。
これは Poisson 射そのものの不存在ではない。
Brauer--Severi ambient 内の三次族、または実際の可逆層を選ばない偏極類を
許す標的は、引き続き有効な候補である。
\end{remark}
